
The random variable $ g(X)$ is a mapping $ g(X): E \to S$, where $ E = \{ 1,\dots, 2n \} $
and $ S = \{ 0,1 \} $. Let's find the distribution of the random variable $ g(X)$. 

Take a subset $ A \subset S = \{ 0,1 \} $
\begin{align*}
    \P_{g(X)} (A) &= \P(g(X) \in A) \\
		  &= \P(\bigcup_{m\in A \subset \{ 0,1 \} } \{ g(X) = m \} )\\
		  &= \sum_{m\in A \subset \{ 0,1 \} } \P( \{ g(X) = m \} )
\end{align*}
where the final line comes from the fact that the sets $ \{ g(X) = m \} $ are disjoint for different
$ m \in \{ 0,1 \} $. 

If $ m= 0$, then by the definition of the random variable $ g(X)$ we have that $ X \leq n$. 
Let's calculate the probability $ \P( \{ g(X) = 0 \}) $

\begin{align*}
    \P( \{ g(X) = 0 \} )&= \P( X(\omega) \leq n \}\\
			&= \P( \{1 \leq X \leq n\} \cup \{ X < 1 \} )\\
			&= \P( \{ 1 \leq X \leq n \} + \P( \{ X < 1 \} 	)
\end{align*}
where the last line follows since $ \{ 1 \leq X \leq n \} $ and $ \{ X <1 \} $ are disjoint. 
We see that since the random variable $ X$ is a mapping	$ X: \Omega \to E = \{ 1,\dots,n \} 
$ we have that the set $ \{ X < 1 \} $ is empty, and the probability of an empty set is zero.
Thus, the probability $ \P( \{ g(X) = 0 \} )$ reduces to

\begin{align*}
    \P( \{ g(X) = 0 \}) &= \P( \{ 1 \leq X \leq n \})\\
			&= \P( \bigcup_{k=1}^{n} \{ X = k \}\\
			&= \sum_{k=1}^{n} \P( \{ X = k \} ) \\
			&= \sum_{k=1}^{n} \frac{1}{2n} \\
			&= 1/2
\end{align*}
where we used the fact that the sets $ \{ X= k \} $ are disjoint for different $ k$ and that the
random variable $ X$ has a uniform distribution on $ \{ 1,\dots, 2n \} $.

If $ m = 1$, then using a similar argument we see that $ X > n$ and so 

\begin{align*}
    \P ( \{ g(X) = 1 \}	) &= \P( \{ X > n \} )\\
			  &= \P( \bigcup_{k=n+1}^{2n} \{ X= k \}    )\\
			  &= \sum_{k=n+1}^{2n} \P(X =k)\\
			  &= \sum_{k=n+1}^{2n} \frac{1}{2n}\\
			  &= 1/2
\end{align*}

All in all, the probability distribution $ P_{g(X)} $ is

\begin{align*}
    P_{g(X)} (A) = \frac{1}{2} |A|
\end{align*}
where $ |A| $ is the size of the subset $ A$, which can have zero, one or two elements.
